% ============================================================
% Chapter 5: Development
%   Section: News Acquisition: Discovery and Extraction
%     Subsection: Monitoring the Sources
% Included from chapters/05-development/02-news-acquisition/news-acquisition.tex
% ============================================================

\subsection{Monitoring the Sources}
\label{subsec:scraping-observability}

A source that stops delivering does not raise an error: ingestion
simply finds nothing new. Three tools make the state of every source
visible, from cheapest to most thorough, and a fourth measures how
quickly articles arrive.

\paragraph{Request statistics.} Every page wanted is one entry, not
every strategy attempt: its final outcome, the strategy that produced
the article, every strategy tried, the HTTP requests actually sent
(none for a guard refusal, one for trafilatura and BeautifulSoup
together, two once the browser steps in), the purpose and the domain.
The counts are thread-safe, persisted in the lake, and shown on the
\texttt{/scraper} page, where ``which step does the work for this
source'' and ``why did it fail'' can be read per domain.

\paragraph{The source probe.} The statistics only know about pages
someone happened to ask for. The probe (\texttt{POST /scraper/probe}
and the ``Source health'' panel) asks on purpose, for each enabled
source: its feed (both feed steps), its topic section pages (always,
since how many links they add is part of the answer), and a real
extraction of a sample of the links found, split between the two. It
classifies each source as:

\begin{itemize}
    \item \textbf{up}: the feed gives links and at least half of the
    sample extracts;
    \item \textbf{degraded}: articles can be had, but not the normal
    way: a dead feed rescued by section pages, or under half the sample
    extracting;
    \item \textbf{down}: nothing extracts.
\end{itemize}

It also sends one query per language to the search engine and reports
which engines answered, since evidence retrieval depends on them. A
probe costs ten to fifteen requests per source, so it runs only when
asked, on a background thread (a second request while one is running
is refused), behind the storage API key, and keeps its last report in
the lake.

\paragraph{The source check.} The probe answers ``is the source up'';
the check (\texttt{POST /sources/check}, the \texttt{/sources} page)
answers ``does every configured page still give a complete article''.
It runs discovery and then extracts a few links through the same
cascade as ingestion, browser included, stores and analyses nothing,
and includes disabled sources. A source is \textbf{working} when every
sample extracted with a title, author and date, \textbf{partial} when
some failed or one lacked those fields, and \textbf{broken} when
discovery found nothing or no sample extracted.

\paragraph{Time to reception.} The freshness report
(\texttt{GET /scraper/freshness}) joins the three times of
Section~\ref{subsec:scraping-metadata} with the lake's own records and
reports, per article and per source, four delays in hours:
\textbf{seen} (publication to first sighting: the feed's delay and the
system's polling), \textbf{queue} (first sighting to fetch: the
system's backlog), \textbf{reception} (publication to fetch) and
\textbf{available} (publication to the reader view). Only precise times
are used; an article dated only by day is reported in days, apart. A
negative delay (a wrong time zone, a rewritten feed time) is counted
rather than averaged in. Because ingestion runs only when someone asks
for it, ``seen'' currently measures mostly the time until the next
round; queue and available are the pipeline's own.
